\documentclass[a4paper,11pt]{article}
\usepackage{iftex}
\ifPDFTeX
\usepackage[T1]{fontenc}
\usepackage[utf8]{inputenc}
\usepackage{lmodern}
\else
\usepackage{fontspec}
\setmainfont{Latin Modern Roman}
\fi
\usepackage[ngerman]{babel}
\usepackage{amsmath,amssymb,amsthm,mathtools}
\usepackage{booktabs,tabularx,array}
\usepackage[a4paper,top=24mm,bottom=25mm,left=27mm,right=27mm]{geometry}
\usepackage{microtype}
\usepackage[hidelinks]{hyperref}
\newcommand{\Pp}[1]{\mathcal P_+(#1)}
\newcommand{\Sing}{\operatorname{Sing}}
\newtheorem{satz}{Satz}
\newtheorem{lemma}{Hilfssatz}
\setlength{\parindent}{0pt}
\setlength{\parskip}{.55em plus .1em minus .1em}
\setlength{\emergencystretch}{1.2em}
\allowdisplaybreaks[2]
\hypersetup{pdftitle={Band 37 - Zweiprodukt-Rekonstruktion - Lesefassung},
 pdfauthor={Martin Kunze},pdfsubject={Rekonstruktion jeder endlichen Halbgruppe mit zweielementigem Produktquadrat}}
\title{Zwei Produktwerte genügen\\[.5em]
\large Rekonstruktion endlicher Halbgruppen\\Lesefassung zu Band 37}
\author{Martin Kunze}
\date{5. Oktober 2026}
\begin{document}
\maketitle

Eine Halbgruppe kann viele Elemente besitzen und trotzdem nur zwei
verschiedene Produkte liefern. Dann liegt jedes Produkt in derselben
zweielementigen Teilhalbgruppe. Ihre innere Multiplikation entscheidet,
welches Rekonstruktionsargument funktioniert: In vier Fällen genügt die
Größe zweier Fasern; im fünften Fall werden die Klassen gleicher Zeilen
und Spalten zurückgewonnen.

Wir schreiben für eine Halbgruppe \(Q\)
\[
 Q^2=\{xy:x,y\in Q\},\qquad Q^3=\{xyz:x,y,z\in Q\}.
\]
Die Potenzhalbgruppe besteht hier aus allen \emph{nichtleeren} Teilmengen:
\[
 \Pp Q=\{A\subseteq Q:A\ne\varnothing\},\qquad
 AB=\{ab:a\in A,\ b\in B\}.
\]
Die leere Menge wird nicht mitgeführt.

\begin{satz}\label{tw:main}
Sei \(S\) eine endliche Halbgruppe mit \(|S^2|=2\). Dann gilt für jede
Halbgruppe \(T\):
\[
 \Pp S\cong\Pp T\quad\Longrightarrow\quad S\cong T.
\]
\end{satz}

Der Satz verlangt weder ein Nullelement noch eine zusätzliche Bedingung
an \(S^3\). Im Fall eines Nullproduktideals folgt \(S^3=\{0\}\) im
Beweis von selbst. Die folgende Darstellung arbeitet alle fünf Fälle
und das abschließende allgemeine Rekonstruktionslemma aus.

\section{Das Produktideal wird algebraisch erkannt}

Sei \(\Phi:\Pp S\to\Pp T\) ein Isomorphismus und
\(C=S^2=\{a,b\}\). Das Produktquadrat ist stets ein beidseitiges
Ideal: Ein weiteres Produkt mit einem beliebigen Grundelement bleibt
ein Produkt zweier Grundelemente. Insbesondere ist \(C\) eine
Teilhalbgruppe von \(S\).

Jedes Produkt nichtleerer Teilmengen liegt in \(C\). Umgekehrt sind
\(\{a\}\) und \(\{b\}\) Produkte geeigneter Einermengen, und
\(SS=C\). Da eine zweielementige Menge nur drei nichtleere Teilmengen
hat, gilt daher
\begin{equation}\label{tw:square}
 (\Pp S)^2=\Pp C=\bigl\{\{a\},\{b\},C\bigr\}.
\end{equation}
Die linke Seite ist allein aus der Multiplikation der Potenzhalbgruppe
bestimmt. Es wird keine Teilmengengröße als zusätzliche Markierung
benötigt.

Zunächst ist auch \(T\) endlich: Seine Einermengen bilden eine
injektive Kopie von \(T\) in die endliche Menge \(\Pp T\). Aus
\[
 2^{|S|}-1=|\Pp S|=|\Pp T|=2^{|T|}-1
\]
folgt \(|S|=|T|\).

Setze \(D=T^2\). Bei \(|D|=1\) hätte \((\Pp T)^2\) nur ein
Element. Bei \(|D|\ge3\) enthielte dieses Produktquadrat mindestens
drei verschiedene Produkt-Einermengen und zusätzlich das davon
verschiedene Produkt \(TT=D\), also mindestens vier Elemente. Wegen
\eqref{tw:square} sind beide Möglichkeiten ausgeschlossen. Somit ist
\(|D|=2\), und \(\Phi\) schränkt sich zu einem Isomorphismus
\begin{equation}\label{tw:restricted}
 \Phi|_{\Pp C}:\Pp C\longrightarrow\Pp D
\end{equation}
ein. Wir müssen zunächst erkennen, welche innere Multiplikation diese
beiden zweielementigen Ideale haben.

\section{Warum genau fünf Typen auftreten}

Auf einer zweielementigen Menge gibt es bis auf Isomorphie genau fünf
Halbgruppen. Das lässt sich unmittelbar aus der Assoziativität
herleiten, ohne alle sechzehn möglichen Tabellen einzeln aufzulisten.

Sind beide Elemente idempotent, gilt \(a^2=a\), \(b^2=b\). Für das
Paar \((ab,ba)\) bleiben vier Möglichkeiten. Die Paare \((a,b)\)
und \((b,a)\) liefern die Linksnull- beziehungsweise
Rechtsnullhalbgruppe. Die Paare \((a,a)\) und \((b,b)\) liefern
denselben Halbverband nach Vertauschung der beiden Namen.

Gibt es genau ein idempotentes Element, können wir es \(a\) nennen.
Dann ist \(a^2=a\), \(b^2=a\). Die Gleichung
\[
 (ba)b=b(ab)
\]
verbietet \(ab\ne ba\): Für \((ab,ba)=(a,b)\) wären die beiden
Seiten \(a\) und \(b\); für \((ab,ba)=(b,a)\) wären sie \(b\)
und \(a\). Also ist \(ab=ba\). Sind beide Werte \(a\), erhalten
wir die Nullhalbgruppe; sind beide Werte \(b\), erhalten wir die
zweielementige Gruppe.

Ohne idempotentes Element müsste \(a^2=b\), \(b^2=a\) gelten.
Dann wäre \((aa)b=bb=a\), während \(a(ab)=b\) gilt, gleichgültig
ob \(ab=a\) oder \(ab=b\). Das widerspricht der Assoziativität.
Die aufgeführten fünf Typen sind selbst assoziativ und erschöpfen
damit alle Möglichkeiten.

\begin{center}
\small
\begin{tabularx}{\linewidth}{@{}l X X@{}}
\toprule
Typ von \(C\) & Multiplikation & Typ von \(\Pp C\) \\
\midrule
Nullhalbgruppe & Alle Produkte sind \(a\).
 & Dreielementige Nullhalbgruppe \\
Linksnullhalbgruppe & \(xy=x\)
 & Dreielementige Linksnullhalbgruppe \\
Rechtsnullhalbgruppe & \(xy=y\)
 & Dreielementige Rechtsnullhalbgruppe \\
Halbverband & \(a^2=ab=ba=a\), \(b^2=b\)
 & Dreielementige Kette \\
Gruppe \(C_2\) & \(a^2=b^2=a\), \(ab=ba=b\)
 & \(C_2\) mit angefügter Null \\
\bottomrule
\end{tabularx}
\end{center}

Die letzte Spalte lässt sich durch die drei Mengen aus
\eqref{tw:square} direkt ausrechnen. Die fünf dort genannten
Halbgruppen sind paarweise nicht isomorph. Links- und Rechtsnulltypen
sind nicht kommutativ und durch die jeweils geltende Identität
\(XY=X\) beziehungsweise \(XY=Y\) verschieden. Die anderen drei
Typen sind kommutativ und besitzen genau ein, drei beziehungsweise
zwei idempotente Elemente.

Im Halbverbandsfall ist die algebraisch bestimmte Ordnung
\(X\le Y\iff XY=X\) genau
\[
 \{a\}<C<\{b\}.
\]
Im Gruppenfall bilden \(\{a\},\{b\}\) die zweielementige Gruppe;
\(C\) ist ein zusätzliches absorbierendes Element. Wegen
\eqref{tw:restricted} müssen \(C\) und \(D\) also denselben Typ
besitzen.

\section{Vier Typen liefern eine Retraktion}

Ist \(C\) keine Nullhalbgruppe, lässt sich jedes Grundelement einem
der beiden Produktwerte zuordnen. Diese Zuordnung \(g:S\to C\)
fixiert die Elemente von \(C\) und bestimmt die ganze Multiplikation:
\begin{equation}\label{tw:inflation}
 g|_C=\operatorname{id}_C,\qquad xy=g(x)g(y)\quad(x,y\in S).
\end{equation}
Eine solche Abbildung heißt hier Retraktion, weil sie auf ihrem
Zielbereich die Identität ist. Wir beweisen ihre Existenz für die vier
Typen einzeln.

\subsection{Halbverband und Gruppe}

Beide zweielementigen Halbgruppen besitzen ein neutrales Element
\(e\). Dieses Element ist zunächst nur innerhalb von \(C\) neutral.
Da \(C\) ein Ideal ist, liegen aber \(ex\) und \(xe\) für jedes
\(x\in S\) in \(C\). Deshalb gilt
\[
 ex=(ex)e=e(xe)=xe.
\]
Setze \(g(x)=ex=xe\). Für \(c\in C\) ist \(g(c)=c\).
Außerdem liegt \(xy\) in \(C\), sodass
\[
 xy=e(xy)e=(ex)(ye)=g(x)g(y).
\]
Damit ist \eqref{tw:inflation} bewiesen.

\subsection{Linksnull und Rechtsnull}

Sei \(C\) eine Linksnullhalbgruppe. Für \(x\in S\) betrachte
\(\lambda_x:C\to C\), \(c\mapsto xc\). Weil \(x^2\in C\)
ist, gilt
\[
 \lambda_x^2(c)=x^2c=x^2.
\]
Die zweite Iteration ist also konstant. Eine nichtkonstante Abbildung
einer zweielementigen Menge ist bijektiv; dann wäre auch ihre zweite
Iteration bijektiv. Daher ist bereits \(\lambda_x\) konstant.
Nenne ihren Wert \(g(x)\).

Für \(x\in C\) ist dieser Wert \(x\). Für beliebige \(x,y\in S\)
und \(c\in C\) folgt wegen \(xy\in C\) und \(yc\in C\)
\[
 xy=(xy)c=x(yc)=g(x)=g(x)g(y).
\]
Auch hier gilt \eqref{tw:inflation}.

Der Rechtsnullfall verläuft spiegelbildlich. Die Abbildung
\(\rho_x:C\to C\), \(c\mapsto cx\), hat die konstante zweite
Iteration \(c\mapsto cx^2=x^2\) und ist daher selbst konstant.
Mit ihrem Wert \(g(x)\) erhalten wir
\[
 xy=c(xy)=(cx)y=g(y)=g(x)g(y).
\]
Die entsprechende Retraktion existiert ebenso in \(T\).

\section{Wie die Potenzhalbgruppe die Fasergrößen verrät}

Für die gerade gewonnenen Retraktionen schreiben wir \(g_S:S\to C\)
und \(g_T:T\to D\). Die beiden Fasern von \(g_S\) haben die
positiven Größen
\[
 m_a=|g_S^{-1}(a)|,\qquad m_b=|g_S^{-1}(b)|.
\]
Sie sind nicht leer, denn sie enthalten \(a\) beziehungsweise \(b\).
Für eine nichtleere Teilmenge \(A\subseteq S\) setzen wir
\[
 G_S(A)=g_S[A]\in\Pp C.
\]
Entsprechend definieren wir \(G_T\). Die drei Fasern von \(G_S\)
haben genau die Größen
\begin{equation}\label{tw:fibres}
 \begin{aligned}
 |G_S^{-1}(\{a\})|&=2^{m_a}-1,\\
 |G_S^{-1}(\{b\})|&=2^{m_b}-1,\\
 |G_S^{-1}(C)|&=(2^{m_a}-1)(2^{m_b}-1).
 \end{aligned}
\end{equation}
Die ersten beiden Zahlen zählen die nichtleeren Teilmengen einer Faser.
Die dritte zählt die Teilmengen, die beide Fasern treffen: Man wählt
unabhängig auf beiden Seiten eine nichtleere Teilmenge.

\subsection{Die Typen mit neutralem Element}

Ist \(C\) ein Halbverband oder die Gruppe \(C_2\), besitzt
\(\Pp C=(\Pp S)^2\) ein eindeutig bestimmtes neutrales Element
\(E_S=\{e_S\}\). Es wird von \(\Phi\) auf das entsprechende
Element \(E_T\) übertragen. Wegen der Definition der Retraktion gilt
\[
 G_S(A)=AE_S,\qquad G_T(B)=BE_T.
\]
Also ist \(G\) algebraisch erkennbar, und für jedes \(A\) gilt
\[
 \Phi(G_S(A))=G_T(\Phi(A)).
\]

Auch die beiden Produkt-Einermengen sind in diesen Typen erkennbar:
Im Halbverband sind sie die beiden Endpunkte der dreielementigen Kette,
in \(C_2\) mit angefügter Null sind sie die beiden Elemente, die nicht
die angefügte Null sind. Es entsteht daher eine Bijektion
\(\alpha:C\to D\) durch
\[
 \Phi(\{c\})=\{\alpha(c)\}\qquad(c\in C).
\]
Sie ist ein Isomorphismus der Ideale, denn \(\Phi\) erhält die
Produkte ihrer Einermengen. Nach \eqref{tw:fibres} gilt für jedes
\(c\in C\)
\[
 2^{|g_S^{-1}(c)|}-1
   =2^{|g_T^{-1}(\alpha(c))|}-1.
\]
Somit stimmen die entsprechend zugeordneten Fasergrößen überein.

\subsection{Die Linksnull- und Rechtsnulltypen}

Hier gilt unmittelbar
\[
 G_S(A)=A^2.
\]
Denn das Produkt zweier Elemente ist im Linksnullfall der
Retraktionswert des ersten, im Rechtsnullfall der Retraktionswert des
zweiten Faktors. Der Isomorphismus \(\Phi\) erhält Quadrate und damit
das Multiset der drei Fasergrößen von \(G\).

Schreibe \(r=2^{m_a}-1\), \(s=2^{m_b}-1\). Dieses Multiset ist
\(\{r,s,rs\}\). Weil \(r,s\ge1\), ist \(rs\) mindestens so groß
wie jede der beiden anderen Zahlen. Die zwei kleinsten Einträge,
\emph{mit ihren Wiederholungen}, sind daher genau \(r,s\). Das gilt
auch bei \(r=1\), \(s=1\) oder \(r=s\). Somit wird das ungeordnete
Paar \(\{m_a,m_b\}\) bestimmt.

Jede Bijektion zwischen zweielementigen Linksnullhalbgruppen ist ein
Isomorphismus; dasselbe gilt für Rechtsnullhalbgruppen. Wir können also
einen Ideal-Isomorphismus \(\alpha:C\to D\) so wählen, dass die
entsprechenden Fasergrößen gleich sind.

\subsection{Aus den Fasern entsteht der Grundisomorphismus}

In allen vier Nichtnullfällen besitzen wir jetzt einen Isomorphismus
\(\alpha:C\to D\) mit
\[
 |g_S^{-1}(c)|=|g_T^{-1}(\alpha(c))|\qquad(c\in C).
\]
Wähle in jeder Faser eine Bijektion, die den ausgezeichneten Punkt
\(c\) auf \(\alpha(c)\) abbildet. Zusammengenommen erhalten wir
eine Bijektion \(f:S\to T\) mit
\[
 f|_C=\alpha,\qquad g_T(f(x))=\alpha(g_S(x)).
\]
Weil \(xy\in C\) liegt und \eqref{tw:inflation} auf beiden Seiten
gilt, folgt
\[
 \begin{aligned}
 f(xy)&=\alpha(g_S(x)g_S(y))\\
      &=g_T(f(x))g_T(f(y))=f(x)f(y).
 \end{aligned}
\]
Damit ist der Hauptsatz für diese vier Typen bewiesen.

\section{Beim Nullideal verschwinden Dreifachprodukte}

Es bleibt der Fall, dass \(C=S^2\) eine Nullhalbgruppe ist. Nenne
ihren konstanten Produktwert \(0\), das andere Element \(u\).
Dann gilt
\[
 S^4=S^2S^2=\{0\},\qquad S^4\subseteq S^3\subseteq S^2=\{0,u\}.
\]
Die Inklusionen folgen durch Zusammenfassen zweier benachbarter
Faktoren. Also kann \(S^3\) nur \(\{0\}\) oder \(S^2\) sein.
Die zweite Möglichkeit ergäbe aber
\[
 S^4=S^3S=S^2S=S^3=S^2,
\]
im Widerspruch zu \(S^4=\{0\}\). Daher folgt automatisch
\begin{equation}\label{tw:nilpotent}
 S^3=\{0\}.
\end{equation}
Weil \(0\in S^2\) ist, liegen \(x0\) und \(0x\) für jedes
\(x\in S\) in \(S^3\). Somit ist \(0\) eine Null der ganzen
Halbgruppe. Für \(T\) gilt dasselbe, denn sein Produktideal hat
denselben Typ. Wir unterscheiden bei Bedarf \(0_S,u_S\) und
\(0_T,u_T\).

\subsection{Die Produkt-Einermengen werden erkannt}

Die eindeutig bestimmte Null der Potenzhalbgruppe ist \(\{0\}\).
Deshalb gilt \(\Phi(\{0_S\})=\{0_T\}\). Um auch die zweite
Produkt-Einermenge zu erkennen, zählen wir Lösungen mit festem linken
Faktor.

Für jede nichtleere Menge \(A\subseteq S\) sind die Lösungen von
\(AB=\{u_S\}\) genau die nichtleeren Teilmengen von
\[
 M_A=\{y\in S:\forall x\in A\ (xy=u_S)\}.
\]
Ihre Zahl ist daher \(2^{|M_A|}-1\). Diese Zahl ist entweder null
oder ungerade.

Der andere Produktwert \(\{0_T,u_T\}\) besitzt dagegen für einen
geeigneten linken Faktor eine \emph{positive gerade} Lösungszahl.
Wähle \(xy=u_T\) und setze \(A=\{0_T,x\}\). Dann löst
\(B=\{y\}\) die Gleichung \(AB=\{0_T,u_T\}\), sodass es
mindestens eine Lösung gibt.

Bei jeder Lösung \(B\) kann man nun die Zugehörigkeit von \(0_T\)
umkehren. Das neue \(B\) bleibt nichtleer, denn wegen des Produktwerts
\(u_T\) muss \(B\) mindestens ein Element verschieden von \(0_T\)
enthalten. Sein Produkt mit \(A\) bleibt gleich: Der linke Faktor
liefert die Null ohnehin, und ein zusätzlicher rechter Nullfaktor
liefert nur nochmals Null. Diese Operation ist eine fixpunktfreie
Involution auf der Lösungsmenge. Daher besteht sie aus Paaren und hat
positive gerade Größe.

Diese Lösungszahlen werden durch \(\Phi\) für jeden festen linken
Faktor übertragen. Also kann \(\Phi(\{u_S\})\) nicht
\(\{0_T,u_T\}\) sein. Es kann wegen der Injektivität auch nicht
\(\{0_T\}\) sein. Da das Zielproduktquadrat nur drei Elemente hat,
folgt
\[
 \Phi(\{u_S\})=\{u_T\}.
\]
Damit werden beide Produkt-Einermengen erhalten. Das folgende allgemeine
Lemma schließt die Rekonstruktion ab.

\section{Ein allgemeines Rekonstruktionslemma}

\begin{lemma}\label{tw:singleton-lemma}
Seien \(S,T\) endliche nichtleere Halbgruppen und
\(\Phi:\Pp S\to\Pp T\) ein Isomorphismus. Gilt
\[
 \Phi[\Sing(S^2)]=\Sing(T^2),\qquad
 \Sing(I)=\{\{c\}:c\in I\},
\]
so ist \(S\cong T\).
\end{lemma}

\begin{proof}
Die Annahme definiert eine Bijektion \(h:S^2\to T^2\) durch
\[
 \Phi(\{c\})=\{h(c)\}\qquad(c\in S^2).
\]
Ferner ist \(|S|=|T|=:n\), weil ihre Potenzhalbgruppen jeweils
\(2^n-1\) Elemente besitzen.

\emph{Erster Schritt: Bilder von Einermengen haben gleiche Zeilen und
Spalten.} Fixiere \(s\in S\) und setze \(A_s=\Phi(\{s\})\).
Für jeden \(c\in S^2\) sei
\[
 k_c=|\{x\in S:sx=c\}|.
\]
Die Gleichung \(\{s\}B=\{c\}\) hat genau \(2^{k_c}-1\)
Lösungen. Ihre Bilder sind alle Lösungen von
\(A_sB=\{h(c)\}\), nämlich die nichtleeren Teilmengen von
\[
 K_c=\{t\in T:\forall a\in A_s\ (at=h(c))\}.
\]
Also gilt \(|K_c|=k_c\). Da \(A_s\) nichtleer und \(h\)
injektiv ist, sind die Mengen \(K_c\) paarweise disjunkt. Ihre
Größen summieren sich zu
\[
 \sum_{c\in S^2}|K_c|=\sum_{c\in S^2}k_c=|S|=|T|.
\]
Sie überdecken daher \(T\). Für jede Spaltenposition \(t\in T\)
besitzen alle Elemente von \(A_s\) denselben Eintrag \(at\).
Also haben sie dieselbe vollständige Zeile.

Die gleiche Rechnung für die Gleichungen \(B\{s\}=\{c\}\)
zeigt, dass alle Elemente von \(A_s\) auch dieselbe vollständige
Spalte haben. Insbesondere liegt \(A_s\) in einer einzigen Klasse
gleicher Zeilen und Spalten in \(T\).

\emph{Zweiter Schritt: Die Klassen und ihre Größen werden übertragen.}
Definiere auf einer Grundhalbgruppe die Äquivalenzrelation
\[
 s\sim t\quad\Longleftrightarrow\quad
 \forall x\ (sx=tx\ \text{und}\ xs=xt).
\]
Auf ihrer Potenzhalbgruppe definieren wir entsprechend
\[
 A\approx B\quad\Longleftrightarrow\quad
 \forall X\ (AX=BX\ \text{und}\ XA=XB).
\]
Die zweite Relation ist allein durch die Multiplikation bestimmt und
wird von \(\Phi\) erhalten.

Sei \(E=[s]_\sim\). Die \emph{vollständige} \(\approx\)-Klasse
von \(\{s\}\) ist genau \(\Pp E\). Tatsächlich erzwingen die
Tests gegen alle Einermengen \(\{x\}\), dass jedes Element einer
äquivalenten Menge dieselbe Zeile und Spalte wie \(s\) hat.
Umgekehrt liefern gleiche Zeilen und Spalten dieselben Produkte auch
gegen beliebige nichtleere Testmengen.

Nach dem ersten Schritt ist \(A_s\) eine nichtleere Teilmenge einer
einzigen Grundklasse \(F\) in \(T\). Deshalb ist seine vollständige
\(\approx\)-Klasse gleich \(\Pp F\). Es folgt
\begin{equation}\label{tw:class-match}
 \Phi(\Pp E)=\Pp F,\qquad
 2^{|E|}-1=2^{|F|}-1,
\end{equation}
also \(|E|=|F|\).

Damit werden verschiedene Grundklassen verschiedenen Grundklassen
zugeordnet. Diese Zuordnung ist auch surjektiv: Die zugeordneten
Klassen sind disjunkt und besitzen zusammen genau \(|S|=|T|\)
Elemente. Es bleibt also keine Zielklasse übrig.

\emph{Dritter Schritt: Die markierten Produktwerte passen in diese
Klassen.} Für jede Zuordnung \(E\leftrightarrow F\) gilt
\begin{equation}\label{tw:marked-match}
 h(E\cap S^2)=F\cap T^2.
\end{equation}
Denn für \(c\in S^2\) ist
\[
 \begin{aligned}
 c\in E
 &\iff\{c\}\approx\{s\}\\
 &\iff\{h(c)\}\approx A_s
 \iff h(c)\in F.
 \end{aligned}
\]
Zusammen mit der Bijektivität von \(h:S^2\to T^2\) ergibt das
\eqref{tw:marked-match}. In jeder Paarung gleich großer endlicher
Klassen kann man daher die vorgeschriebene Teilbijektion \(h\) zu
einer ganzen Klassenbijektion erweitern. Zusammengenommen erhalten wir
eine Bijektion \(f:S\to T\) mit \(f|_{S^2}=h\).

\emph{Vierter Schritt: Die Bijektion erhält die Multiplikation.}
Alle Elemente von \(A_s\) haben dieselbe Zeile und Spalte wie
\(f(s)\). Für \(s,t\in S\) folgt daher
\[
 A_sA_t=\{f(s)f(t)\}.
\]
Hier müssen \(f(s)\) und \(f(t)\) selbst nicht in \(A_s,A_t\)
liegen; ihre gleichen Zeilen und Spalten genügen. Andererseits gilt
\[
 A_sA_t=\Phi(\{s\}\{t\})=\Phi(\{st\})
       =\{h(st)\}=\{f(st)\}.
\]
Damit ist \(f(s)f(t)=f(st)\). Die Bijektion \(f\) ist also ein
Halbgruppenisomorphismus.
\end{proof}

Im verbleibenden Nullidealfall wurden die Voraussetzungen dieses Lemmas
im vorigen Abschnitt nachgewiesen. Es liefert daher auch dort
\(S\cong T\). Zusammen mit den vier durch Retraktionen behandelten
Typen ist Satz~\ref{tw:main} vollständig bewiesen.\qed

\clearpage
\section*{Anschluss an das Skript}

Die bisherige Lesefassung \emph{Was Nullprodukte verraten} behandelt
endliche Halbgruppen mit einem Nullelement und genau zwei Produktwerten.
Der hier bewiesene Satz umfasst alle endlichen Halbgruppen mit
zweielementigem Produktquadrat. Dafür muss dessen innere Multiplikation
unterschieden werden. Die Nullhalbgruppen- und Halbverbandstypen haben
auch in der Grundhalbgruppe eine Null und gehören zum früheren
Zweiproduktsatz: Im Halbverbandstyp folgt dies aus
\eqref{tw:inflation}, im Nullidealfall aus \eqref{tw:nilpotent}.
Der vorliegende Satz umfasst zusätzlich beide einseitigen Nulltypen
und den Gruppentyp. Im Nullidealfall folgen die verschwindenden
Dreifachprodukte aus der Struktur des Produktideals selbst.

Die allgemeinen Grundlagen gehören zu Band~20 (endliche Zählung),
Band~28 (Halbgruppen und Potenzhalbgruppen), Band~33 (Nullproduktprofile)
und Band~37 (endliche Halbgruppen). Die Lesefassung ist eigenständig:
Alle für ihren Hauptsatz benötigten besonderen Argumente stehen oben.

Ausgangspunkt ist der am 5.~Oktober 2026 bereitgestellte Beweisentwurf.
Die mathematische Gegenprüfung wurde in die Darstellung eingearbeitet,
insbesondere die direkte Klassifikation und die Vereinbarkeit der
markierten Produktwerte mit den Klassenbijektionen. Eine Aussage zur
literaturgeschichtlichen Priorität wird damit nicht verbunden. Der Satz
behandelt das zweielementige Produktquadrat; die allgemeine endliche
Potenzhalbgruppenvermutung umfasst weitere Fälle.

\end{document}
